% --- IMPORTS ---
\documentclass[leqno]{article}
\usepackage{graphicx}
\usepackage{dsfont}
\usepackage{mathtools}
\usepackage{amssymb}
\usepackage{amsmath}
\usepackage{amsthm}
\usepackage{float}
\usepackage{ragged2e}
\usepackage{tensor}
\usepackage{fancyhdr}
\usepackage{empheq}
\usepackage[most]{tcolorbox}
\usepackage{enumitem}
\usepackage{dirtytalk}
\usepackage{marginnote}
\usepackage{microtype}
\usepackage[
  a4paper,
  left=0.8in,
  right=1in,
  top=1in,
  bottom=0.5in,
  headheight=14pt,
  headsep=0.4in
]{geometry}
\setlength{\parindent}{0cm}
% --- TikZ Packages ---
\usepackage{pgfplots}
\pgfplotsset{compat=1.18}
\usepgfplotslibrary{fillbetween, polar}
\usetikzlibrary{calc, positioning}
\usetikzlibrary{angles, quotes}
\usetikzlibrary{arrows.meta}
\usetikzlibrary{decorations.pathreplacing, decorations.markings}
\usetikzlibrary{patterns}
\usetikzlibrary{intersections}
% --- URL Packages ---
\usepackage[colorlinks=true,linkcolor=red,citecolor=magenta,urlcolor=black]{hyperref}%
\urlstyle{same}
% --- END OF IMPORTS ---

% --- CUSTOM COMMANDS ---
\RenewDocumentCommand{\marks}{o m}{%
  \IfValueTF{#1}%
    {\marginnote{\raggedright\textbf{[#2]}}[#1]}%
    {\marginnote{\raggedright\textbf{[#2]}}}%
}
\newcommand{\vspacerule}[2]{%
  \par\vspace*{#1}%
  \noindent\hrulefill\par
  \vspace*{#2}%
}
\definecolor{scarlet}{RGB}{205, 40, 75}
\newcommand{\questionitem}{%
  \item\phantomsection
  \addcontentsline{toc}{section}{Question \number\value{enumi}}%
}
\renewcommand{\contentsname}{Questions}
% --- END OF CUSTOM COMMANDS ---

% --- HEADER CONFIG ---
\pagestyle{fancy}
\fancyhead{}
\fancyfoot{}
\renewcommand{\headrulewidth}{0.9pt}
\lhead{Vectors - Reflections in Planes}
\chead{}
\rhead{\href{https://danieljsmith.org}{danieljsmith.org}}
% --- END OF HEADER CONFIG ---

\begin{document}
\vspace*{20pt}
\tableofcontents
\newpage
\begin{enumerate}[
  label=\textbf{\arabic*.},
  leftmargin=*,
  topsep=0pt,
  partopsep=0pt,
  parsep=0pt
]


% ---- Question 1 ----
\questionitem
% [Source: _QBT__Reflections_in_Planes, Q1]

The line $\ell_1$ is the intersection of the planes
\[
x+y+z=1
\] \\
and
\[
2x-y+z=4
\] \\
The plane $\Pi$ has equation
\[
x+y+2z=3
\] \\
The line $\ell_2$ is the reflection of the line $\ell_1$ in the plane $\Pi$. \\

Find a vector equation of the line $\ell_2$. \marks{7}

\newpage

% ---- Question 2 ----
\questionitem
% [Source: _QBT__Reflections_in_Planes, Q2]

The plane $\Pi$ has vector equation
\[
\mathbf{r}=\begin{pmatrix}1\\2\\-1\end{pmatrix}+s\begin{pmatrix}2\\-1\\1\end{pmatrix}+t\begin{pmatrix}1\\1\\0\end{pmatrix}
\] \\
The line $\ell_1$ passes through the point $P(4,1,2)$ and is perpendicular to $\Pi$. \\
The line $\ell_1$ meets $\Pi$ at the point $Q$. \\
\begin{enumerate}[label=\textbf{(\alph*)}]
  \item Find the coordinates of $Q$. \marks{4} \\
\end{enumerate}
Given that the point $R(3,1,0)$ lies on $\Pi$, \\
\begin{enumerate}[label=\textbf{(\alph*)}, resume]
  \item find a Cartesian equation of the plane containing $P$, $Q$ and $R$. \marks{4} \\
\end{enumerate}
The line $\ell_2$ passes through $P$ and $R$. \\
The line $\ell_3$ is the reflection of $\ell_2$ in $\Pi$. \\
\begin{enumerate}[label=\textbf{(\alph*)}, resume]
  \item Find a vector equation of $\ell_3$. \marks{4} \\
\end{enumerate}

\newpage

% ---- Question 3 ----
\questionitem
% [Source: _QBT__Reflections_in_Planes, Q3]

A plane passes through the points $C(1,-1,2)$, $D(3,0,1)$ and $E(2,2,2)$. \\

The point $A$ has coordinates $(4,-2,7)$ and the point $B$ is the reflection of $A$ in the plane. \\

Find the coordinates of the point $B$. \marks{7}

\newpage

% ---- Question 4 ----
\questionitem
% [Source: _QBT__Reflections_in_Planes, Q4]

A line $\ell$ has equation
\[
\mathbf{r}=\begin{pmatrix}4\\2\\-3\end{pmatrix}+\lambda\begin{pmatrix}2\\-1\\1\end{pmatrix}
\] \\
Another line $m$ has equation
\[
\mathbf{r}=\begin{pmatrix}1\\1\\2\end{pmatrix}+\mu\begin{pmatrix}2\\1\\0\end{pmatrix}
\] \\
A plane $\Pi$ contains $m$ and is parallel to $\ell$. \\
\begin{enumerate}[label=\textbf{(\roman*)}]
  \item Find an equation of $\Pi$, giving your answer in the form $\mathbf{r}\cdot\mathbf{n}=p$. \marks{4} \\
  \item Find the distance between $\ell$ and $\Pi$. \marks{4} \\
  \item Find an equation of the line which is the reflection of $\ell$ in $\Pi$, giving your answer in the form $\mathbf{r}=\mathbf{a}+t\mathbf{b}$. \marks{4} \\
\end{enumerate}

\newpage

% ---- Question 5 ----
\questionitem
% [Source: _QBT__Reflections_in_Planes, Q5]

The line $\ell_1$ has vector equation
\[
\mathbf{r}=\begin{pmatrix}1\\4\\0\end{pmatrix}+\lambda\begin{pmatrix}2\\1\\1\end{pmatrix}
\] \\
The plane $\Pi_1$ passes through the point $(2,0,0)$ and has normal vector
\[
\begin{pmatrix}3\\-1\\2\end{pmatrix}
\] \\
\begin{enumerate}[label=\textbf{(\alph*)}]
  \item Find the point of intersection of $\ell_1$ and $\Pi_1$. \marks{2} \\
\end{enumerate}
The line $\ell_2$ is the reflection of the line $\ell_1$ in the plane $\Pi_1$. \\
\begin{enumerate}[label=\textbf{(\alph*)}, resume]
  \item Show that a vector equation for the line $\ell_2$ is
  \[
  \mathbf{r}=\begin{pmatrix}4\\3\\2\end{pmatrix}+\mu\begin{pmatrix}-1\\2\\-1\end{pmatrix}
  \]
  where $\mu$ is a scalar parameter. \marks{5}\\
\end{enumerate}
The plane $\Pi_2$ contains the line $\ell_1$ and the line $\ell_2$. \\
\begin{enumerate}[label=\textbf{(\alph*)}, resume]
  \item Determine a vector equation for the line of intersection of $\Pi_1$ and $\Pi_2$. \marks{2} \\
\end{enumerate}

\newpage

% ---- Question 6 ----
\questionitem
% [Source: _QBT__Reflections_in_Planes, Q6]

The triangle $ABC$ lies in the plane $\Pi$. The vertices are
\[
A(1,0,0),\ B(2,2,1) \text{ and } C(1,1,c)
\] \\
where $c$ is a constant. \\
\begin{enumerate}[label=\textbf{(\alph*)}]
  \item Find, in terms of $c$, $\overrightarrow{BA}\times\overrightarrow{BC}$. \marks{3} \\
\end{enumerate}
Given that $\overrightarrow{BA}\times\overrightarrow{BC}=d\mathbf{i}+2\mathbf{j}-\mathbf{k}$, where $d$ is a constant, \\
\begin{enumerate}[label=\textbf{(\alph*)}, resume]
  \item find the value of $c$ and show that $d=-3$. \marks{2} \\
  \item find an equation of $\Pi$ in the form $\mathbf{r}\cdot\mathbf{n}=p$, where $p$ is a constant. \marks{3} \\
\end{enumerate}
The point $S$ has position vector $3\mathbf{i}+\mathbf{j}+2\mathbf{k}$. The point $S'$ is the image of $S$ under reflection in $\Pi$. \\
\begin{enumerate}[label=\textbf{(\alph*)}, resume]
  \item Find the position vector of $S'$. \marks{5} \\
\end{enumerate}

\newpage

% ---- Question 7 ----
\questionitem
% [Source: _QBT__Reflections_in_Planes, Q7]

The line $\ell$ has equation
\[
\mathbf{r}=(\mathbf{i}+\mathbf{k})+\lambda(\mathbf{i}+2\mathbf{j}+3\mathbf{k})
\] \\
where $\lambda$ is a scalar parameter, and the plane $\Pi$ has equation
\[
\mathbf{r}\cdot(2\mathbf{i}-\mathbf{j}+\mathbf{k})=9
\] \\
\begin{enumerate}[label=\textbf{(\alph*)}]
  \item Find the coordinates of the point of intersection of $\ell$ and $\Pi$. \marks{4} \\
\end{enumerate}
The perpendicular to $\Pi$ from the point $A\,(0,4,1)$ meets $\Pi$ at the point $B$. \\
\begin{enumerate}[label=\textbf{(\alph*)}, resume]
  \item Verify that the coordinates of $B$ are $(4,2,3)$. \marks{3} \\
\end{enumerate}
The point $A\,(0,4,1)$ is reflected in the plane $\Pi$ to give the image point $A'$. \\
\begin{enumerate}[label=\textbf{(\alph*)}, resume]
  \item Find the coordinates of the point $A'$. \marks{2} \\
  \item Find an equation for the line obtained by reflecting the line $\ell$ in the plane $\Pi$, giving your answer in the form
  \[
  \mathbf{r}\times\mathbf{a}=\mathbf{b}
  \]
  \marks[-20pt]{4}
\end{enumerate}


\end{enumerate}
\end{document}
